% Plan: development/outline.md, section 6, item F. Proof of the padding lemma
% from development/eg-lemma-A1.md, whose underflow argument is replaced here;
% audit code and results in development/eg-audit/ (out/compare.log,
% out/run_audit.log, out/logs/); corrections and the trust-base wording from
% the separate review development/reviews/sol-eg-audit.md (items 1-4) and
% development/revision-notes.md.
\section{Rounding-error analysis for the \inst{eg} certificates}\label{app:egrounding}

The displayed dual bounds of \cref{thm:eg-bounds} rest on the second code, the leaf re-certifier (\cref{app:eg-computation}), which pads its binary64 enclosures explicitly.
We prove that the paddings suffice if every exponential and integer power that the program obtains from NumPy is accurate to a relative $10^{-14}$ (\cref{lem:eg-padding}).
On the machine used, NumPy~2.5.1 evaluates these functions with Intel SVML routines (through its AVX-512 loops), for which we know no proved error bound, so an accuracy auditor checked every such value used in a rerun on all leaves (\cref{app:egrounding-auditor,app:egrounding-run}); \cref{thm:eg-trust} states the resulting trust base.

\subsection{Setting, hypotheses and statement}\label{app:egrounding-setting}

The re-certifier reads the data of each instance from the MINLPLib GAMS file as exact rationals; an exact comparison found them identical to the OSIL data.
In this section a star marks the exact data and the exact quantities of \cref{app:eg-lemmas}: $a^*_m$, $z^*_{mi}$, $\gamma^*_i<0$, $s^*_i\in\{1,\frac1{10}\}$, $\lambda^*_i$, and $t^*_{mi}$, $E^*_m$, $w^*_m$, $S_0^*$, $S_1^*$, $S_2^*$, $T^*$, $K_1^*$, $K_D^*$, $G^*$ (for $G_0$), $\nabla^*$, $H^*$, $\ell^*_m$, $q^*$, $P_2^*$, $P_3^*$, formed with the exact data at the float centre $c$.
The exact row function is $g^*(x)=\sum_ma^*_m\exp\bigl(\sum_i\gamma^*_i(s^*_ix_i-z^*_{mi})^2\bigr)+\lambda^*\cdot x$.
The program rounds the term data to doubles, $\hat a=\fl(a^*)$, $\hat z=\fl(z^*)$, $\hat\gamma=\fl(\gamma^*)$, $\hat s=\fl(s^*)$ and $\hat\lambda=\fl(\lambda^*)$; it keeps $c_k$, the side-row bounds, the variable bounds and the target $\theta^*_I$ exact for the final tests.
Symbols without a star denote floats computed by the program.
The tables below name some of them as in the archived code, to make each line easy to check against it.

For these programs we use \cref{hyp:fp-ieee} in the following form.
NumPy float64 arithmetic is IEEE~754 binary64 with round-to-nearest-even and gradual underflow; comparisons, \texttt{abs}, \texttt{max}, \texttt{min}, \texttt{where}, \texttt{clip}, \texttt{floor}, \texttt{rint} and int64 integer operations are exact; Python integer and \texttt{Fraction} arithmetic is exact, and the conversion of a rational to a double is correctly rounded.
Sums and dot products may be evaluated in any order, with or without fused multiply-add.

Let $\varepsilon=10^{-14}$.
The two conditions below concern the results $y$ of the calls that the enclosure procedures make for a box $X=[lo,hi]$ (floats, $lo\le hi$, integer coordinates relaxed to intervals):
\begin{align}
  &\text{every call } y=\exp(x)\text{ satisfies } |y-e^x|\le\varepsilon e^x \text{ if } x\ge-708, \text{ and } 0\le y\le2^{-1020} \text{ if } x<-708; \tag{E}\label{eq:eg-E}\\
  &\text{every call } y=x^k\ (k\in\{2,3,4\}) \text{ satisfies } x\ge0 \text{ and } |y-x^k|\le\varepsilon x^k. \tag{P}\label{eq:eg-P}
\end{align}
The root boxes of the search contain the box of the model, slightly enlarged outward where its ends are not floats; every leaf and every piece of a split leaf lies in a root box.

\begin{lemma}[paddings of the leaf re-certifier]\label{lem:eg-padding}
Assume \cref{hyp:fp-ieee} in the form above.
Let $X=[lo,hi]$ be a box of floats contained in a root box of the search for an instance $I$, and assume \eqref{eq:eg-E} and \eqref{eq:eg-P} for the calls made for $X$.
Let the natural enclosure return $n^{\mathrm{lo}}_k,n^{\mathrm{hi}}_k$ and the Taylor model return the centre $c$, the gradients $\beta_k$ and the constants $aL_k,aU_k$, with the program's assertion $\delta E<10^{-6}$ satisfied.
Then for every row $k$ and every $x\in X$:
\begin{enumerate}
\item[(i)] $n^{\mathrm{lo}}_k\le g^*_k(x)\le n^{\mathrm{hi}}_k$;
\item[(ii)] $aL_k+\beta_k\cdot(x-c)\le g^*_k(x)\le aU_k+\beta_k\cdot(x-c)$.
\end{enumerate}
That is, the computed floats satisfy \eqref{eq:eg-rowbounds} for the exact rows.
\end{lemma}

Infinite values satisfy (i) and (ii) trivially, and the re-certifier ignores them; a NaN would make the exact tests raise an error.

\subsection{Rounding facts}\label{app:egrounding-facts}

Let $\uround=2^{-53}$, $\gamma_n=n\uround/(1-n\uround)$ and $\eta_0=2^{-1075}$.
We use the following consequences of \cref{hyp:fp-ieee}.
\begin{enumerate}
\item[(F0)] \emph{One operation.} If $z$ is the exact result of $+,-,\times,\div$ on floats and $\hat z$ the computed one, then $|\hat z-z|\le\uround|z|$ and $|\hat z-z|\le\uround|\hat z|$, unless the result is subnormal.
A subnormal sum or difference is exact; a subnormal product or quotient has error at most $\eta_0$.
Rounding is monotone: $z_1\le z_2$ implies $\hat z_1\le\hat z_2$.
\item[(F1)] \emph{Data.} The double nearest to a rational differs from it by a relative $\uround$ at most.
$\hat s=1$ is exact, and $\hat s=\fl(0.1)$ exceeds $0.1$ by a relative $\frac12\uround$; so $\hat s\ge s^*$ and $\hat s^2\ge s^{*2}$.
\item[(F2)] \emph{Sums.} A float sum of $n$ terms, in any order, differs from the exact sum by at most $\gamma_{n-1}\sum|x_i|$: each term passes through at most $n-1$ roundings, and $|\prod_{j\le J}(1+\delta_j)-1|\le\gamma_J$ for $|\delta_j|\le\uround$.
\item[(F3)] \emph{Sums of products.} A float sum of $n$ products of $k$ factors, in any order and with or without fused multiply-add, has error at most $\gamma_{n+k-2}\sum|\text{products}|$, by the same argument.
\item[(F4)] \emph{Computed pads.} A nonnegative expression in nonnegative floats evaluated with $j$ operations $+$ and $\times$ is computed with a value of at least $(1-\uround)^j$ times its exact value.
\item[(F5)] \emph{Directed steps.} If a pad $e\ge0$ is subtracted from a float $v$, then $\fl(v-e)\le v-e+\uround|\fl(v-e)|$; multiplying by a float factor $1+\delta$ is exact up to a relative $\uround$.
A pad applied by a float operation must therefore exceed the error it covers by $\uround$ times the padded value; every line below includes this.
\end{enumerate}
The code's constants written $10^{-15}$, $10^{-14}$, $10^{-13}$, $10^{-12}$, $1.1\cdot10^{-14}$ and $1.02$ are doubles within a relative $8\cdot10^{-17}$ of these decimals.
The factors written $1\mp10^{-14}$, $1+10^{-15}$, $1+10^{-13}$ and $1+10^{-12}$ are the doubles $1\mp\sci{9.992007}{-15}$, $1+\sci{1.110223}{-15}$, $1+\sci{9.992007}{-14}$ and $1+\sci{1.0000889}{-12}$.
All headroom values below were computed with these exact values and are rounded down.
We use $\uround\le\sci{1.111}{-16}$, $\varepsilon=90.07\,\uround$, $\gamma_6\le\sci{6.67}{-16}$, $\gamma_{96}\le\sci{1.066}{-14}$, $\gamma_{97}\le\sci{1.077}{-14}$, $\gamma_{98}\le\sci{1.089}{-14}$, $\gamma_{99}\le\sci{1.100}{-14}$ and $\gamma_{348}\le\sci{3.864}{-14}$.
Every pad is itself evaluated in floats.
By (F4) its computed value is at least $(1-j\uround)$ times its exact value, with $j\le10$ for per-element pads and $j\le100$ for 97-term pads; a headroom of 2 or more absorbs this, and where the headroom is smaller the stated tolerance includes it.

\paragraph{Underflow.}
The tables treat relative errors only; underflow is covered as follows.
We computed bounds from the exact data and the root boxes in rational arithmetic.
On the root boxes the exact exponents satisfy $-14.25\le E\le0$, so the branches of the code for exponents below $-699$ are never taken and no exponential result is subnormal.
The data and root boxes also give $|\hat a|\le1712$, $\sum_m|\hat a_m|\le39{,}667$, $|t|\le3.01$, $|k_{1,i}|r_i\le13.91$, $\ell_m\le16.97$ and $q\le5.14$, so that $e^{\ell_m}<\sci{2.4}{7}$ and $\ell_m+2q\le27.3$.
With these bounds, which hold with room to spare near the computed values, the first-order sensitivity of every output ($n^{\mathrm{lo}}$, $n^{\mathrm{hi}}$, $aL$, $aU$, and the pad $e_\beta$ against $|\beta-\nabla^*|$) to any intermediate result is below $10^{16}$.
The largest factors occur in the fourth-order remainder: its sensitivity to $q$ is at most $\sum_m|\hat a_m|\,e^{16.97}\cdot7{,}401<\sci{7}{15}$, and those to $\ell_m$ and to the exponents $E_{0,m}$ are below $\sci{1.2}{15}$.
An underflow changes one intermediate result by at most $\eta_0$, and fewer than $10^6$ operations feed one output.
So all underflow errors of one output total less than $10^{22}\eta_0<10^{-301}$.
Every output ends with a relative pad, at least $10^{-13}$ times a sum $\sigma$ that bounds the output in magnitude, followed by an absolute pad of $10^{-300}$.
If $\sigma>\sci{4}{-285}$, the part of the relative pad that the rounding errors do not use exceeds $10^{-298}$.
Otherwise the output is at most about $\sci{4}{-285}$ in magnitude, and adding or subtracting $10^{-300}$ moves it by at least $10^{-300}-\uround\cdot\sci{4.02}{-285}>\sci{5}{-301}$.
Either way the underflow errors are absorbed.

\subsection{Proof of \texorpdfstring{\cref{lem:eg-padding}}{the padding lemma}}\label{app:egrounding-proof}

For each computed quantity, \cref{tab:eg-natural,tab:eg-taylor} give the property it must have, the worst-case error that the pad must cover (including the rounding of the padding operations, (F5)), the pad in the code, and the headroom: the exact value of the pad divided by the error it covers, or the largest relative error of the exponential or power that the line tolerates.

% Table cells of the two long tables: small type, ragged right.
\providecommand{\egcell}[1]{\scriptsize\raggedright\let\\\tabularnewline #1}

\paragraph{Natural enclosure.}
For $x_i\in[lo_i,hi_i]$, $t^*_{mi}(x)=s^*_ix_i-z^*_{mi}$ is increasing in $x_i$.
The program computes $t^{\mathrm{l}}=\fl(\fl(\hat s\,lo)-\hat z)$, $t^{\mathrm{h}}=\fl(\fl(\hat s\,hi)-\hat z)$, widens them by $\delta t$, and proceeds as in \cref{lem:eg-natural}.

\begingroup\setlength{\LTcapwidth}{\textwidth}\setlength{\tabcolsep}{3pt}\renewcommand{\arraystretch}{1.0}
\begin{longtable}{@{}p{1.9cm}p{2.8cm}p{5.2cm}p{4.0cm}p{1.25cm}@{}}
\caption{Natural enclosure: errors and pads.}\label{tab:eg-natural}\\
\toprule
\egcell{quantity} & \egcell{property} & \egcell{error to cover} & \egcell{pad} & \egcell{headroom}\\
\midrule
\endfirsthead
\toprule
\egcell{quantity} & \egcell{property} & \egcell{error to cover} & \egcell{pad} & \egcell{headroom}\\
\midrule
\endhead
\bottomrule
\endfoot
\egcell{ends $t^{\mathrm{l}},t^{\mathrm{h}}$ (\texttt{tl}, \texttt{th})} &
\egcell{after widening, $t^{\mathrm{l}}\le t^*_{mi}(lo_i)$ and $t^{\mathrm{h}}\ge t^*_{mi}(hi_i)$} &
\egcell{data $\uround|\hat z|$; product and $\hat s$ vs $s^*$: $2.01\uround|\hat sx|$; sum $\uround|t|$; widening step $\uround(|t|+\delta t)$} &
\egcell{$\delta t=10^{-15}(|\hat z|+\hat s\max(|lo|,|hi|)+\max(|t^{\mathrm{l}}|,|t^{\mathrm{h}}|))+10^{-300}$} &
\egcell{4.48}\\
\egcell{squares $\theta^{\min},\theta^{\max}$ (\texttt{sqmin}, \texttt{sqmax})} &
\egcell{bounds on $t^{*2}$ over the widened interval} &
\egcell{one rounding ($\uround$), carried to the next line} & \egcell{---} & \egcell{---}\\
\egcell{exponent ends $E^{\mathrm{lo}},E^{\mathrm{hi}}$} &
\egcell{$E^{\mathrm{lo}}\le E^*_m(x)\le E^{\mathrm{hi}}$ on $X$} &
\egcell{data $\uround$, square $\uround$, product $\uround$, 7-term sum $\gamma_6$, widening $\uround$: at most $10\uround\sum|\hat\gamma|\theta^{\max}$} &
\egcell{$10^{-14}\sum|\hat\gamma|\theta^{\max}+10^{-300}$} & \egcell{9.0}\\
\egcell{exponentials $\xi^{\mathrm{lo}},\xi^{\mathrm{hi}}$ (\texttt{elo}, \texttt{ehi})} &
\egcell{approximations of $e^{E^{\mathrm{lo}}}$ and $e^{E^{\mathrm{hi}}}$; they need not bound them, because their errors are covered in the next line} &
\egcell{\eqref{eq:eg-E}; one product; $\xi^{\mathrm{lo}}=0$ if $E^{\mathrm{lo}}<-700$} &
\egcell{factors $1\mp10^{-14}$; $+10^{-300}$ on $\xi^{\mathrm{hi}}$} & \egcell{chained to the next line}\\
\egcell{term ends $b^{\mathrm{lo}}_m,b^{\mathrm{hi}}_m$ (\texttt{tlo}, \texttt{thi})} &
\egcell{$b^{\mathrm{lo}}_m\le a^*_me^{E^*_m(x)}\le b^{\mathrm{hi}}_m$ on $X$} &
\egcell{for $\hat a>0$, lower end: $(1+\varepsilon)(1-\sci{9.992}{-15})(1+\uround)^4(1-10^{-14}(1-\uround))\le1$ is needed; for $\hat a<0$ the mirror condition through $\xi^{\mathrm{hi}}$} &
\egcell{$\mp10^{-14}|\cdot|\mp10^{-300}$} &
\egcell{holds for $\varepsilon\le\sci{1.954}{-14}$}\\
\egcell{97-term sums} &
\egcell{lower and upper bounds on $\sum_m b_m$} &
\egcell{$\gamma_{96}\sum|b_m|$ and two pad steps} &
\egcell{$10^{-13}\sum|b_m|+10^{-300}$} & \egcell{9.1}\\
\egcell{linear part, outputs $n^{\mathrm{lo}},n^{\mathrm{hi}}$} &
\egcell{add $\min$ and $\max$ of $\lambda^*\cdot x$ over $X$} &
\egcell{data $\uround$, product $\uround$, 7-term sum $\gamma_6$, final sum $\uround$, pad steps $2\uround$} &
\egcell{$10^{-13}(\sum|\text{linear terms}|+|\text{sum of terms}|)+10^{-300}$} & \egcell{$>10$}\\
\end{longtable}
\endgroup

In the line for the exponent ends, the error of $E^{\mathrm{hi}}$ is bounded by the same expression because $\theta^{\min}\le\theta^{\max}$.
In the line for the term ends, the two conditions hold for $\varepsilon$ up to about $\sci{1.95479}{-14}$.
When $E^{\mathrm{hi}}<-708$, condition \eqref{eq:eg-E} only gives $y\ge0$; then $\xi^{\mathrm{hi}}\ge10^{-300}(1-\uround)>e^{E^{\mathrm{hi}}}$ because of the absolute pad.
When $E^{\mathrm{lo}}<-700$, the lower term end of a positive weight is $-10^{-300}$, below the positive exact term.
Since the exponential is increasing and each term lies between its two ends, item~(i) of \cref{lem:eg-padding} follows.

\paragraph{Taylor model.}
The program forms the quantities of \cref{lem:eg-taylor} at the float centre $c$ from rounded data and pads each of them.
Let $\mathcal D\subseteq\{1,\dots,7\}$ be the coordinates with $r_i>0$ for some box of the batch; for $i\notin\mathcal D$ every box has $r_i=0$, so $d_i=0$ and these coordinates drop out.
If $\mathcal D$ is empty, $X$ is a point and the program sets $\beta=0$ and $P=Q_L=Q_U=0$, which is correct because $d=0$.

\begingroup\setlength{\LTcapwidth}{\textwidth}\setlength{\tabcolsep}{3pt}\renewcommand{\arraystretch}{1.0}
\begin{longtable}{@{}p{1.9cm}p{2.8cm}p{5.2cm}p{4.0cm}p{1.25cm}@{}}
\caption{Taylor model: errors and pads.
$\tau_m=\max_{i\in\mathcal D}\fl(|t_{mi}|+\delta t_{mi})$ and $\delta\tau_m=\max_{i\in\mathcal D}\delta t_{mi}$.}\label{tab:eg-taylor}\\
\toprule
\egcell{quantity} & \egcell{property} & \egcell{error to cover} & \egcell{pad} & \egcell{headroom}\\
\midrule
\endfirsthead
\toprule
\egcell{quantity} & \egcell{property} & \egcell{error to cover} & \egcell{pad} & \egcell{headroom}\\
\midrule
\endhead
\bottomrule
\endfoot
\egcell{centre $c$, radii $r$} &
\egcell{$c\in X$; $r_i\ge\max(hi_i-c_i,c_i-lo_i)$} &
\egcell{$\fl(hi-c)\ge(hi-c)/(1+\uround)$; product $\uround$} &
\egcell{factor $1+10^{-15}$; $c$ clipped to $X$} & \egcell{5.0}\\
\egcell{$t$, $\delta t$} &
\egcell{$|t-t^*|\le\delta t$} &
\egcell{$\uround|\hat z|+2.01\uround|\hat sc|+\uround|t|$} &
\egcell{$10^{-15}(|\hat z|+|\hat sc|+|t|)+10^{-300}$} & \egcell{4.48}\\
\egcell{$E_0$, $\delta E$} &
\egcell{$|E_0-E^*_0|\le\delta E$; the program asserts $\delta E<10^{-6}$} &
\egcell{from $t$: $|\hat\gamma|(1+\uround)(2|t|+e_t)e_t$ with $e_t\le\delta t/4.48$; rounding: data, two products, 7-term sum, at most $9\uround\sum|\hat\gamma|t^2$} &
\egcell{$\sum|\hat\gamma|\bigl[(2|t|+\delta t)\delta t+10^{-14}t^2\bigr]+10^{-300}$} &
\egcell{4.48 (data); 10 (rounding)}\\
\egcell{$w$, $\delta w$} &
\egcell{$|w-w^*|\le\delta w$} &
\egcell{with $w=\fl(\hat a\,y)$, $y=\exp(E_0)$: $|w|\,\bigl(e^{\delta E}/((1-\uround)^2(1-\varepsilon))-1\bigr)\le|w|(1.000001\,\delta E+\varepsilon+2.01\uround)$} &
\egcell{$|w|(1.02\,\delta E+1.1\cdot10^{-14})+10^{-300}$ (and $|\hat a|10^{-300}$ if $E_0<-699$; $w=0$ if $E_0<-700$)} &
\egcell{holds for $\varepsilon\le\sci{1.077}{-14}$}\\
\egcell{$\bar w=|w|+\delta w$ (\texttt{AW})} &
\egcell{$\bar w\ge|w^*|(1-\uround)$} &
\egcell{one rounding} &
\egcell{absorbed by the factors $1+10^{-12}$ on $P$} & \egcell{---}\\
\egcell{$S_0$, $e_0$; $G$, $e_G$} &
\egcell{$|S_0-S_0^*|\le e_0$; $|G-G^*|\le e_G$} &
\egcell{$\sum\delta w+\gamma_{96}\sum|w|$; linear part at $c$ and final sum: at most $9\uround(|S_0|+\sum|\hat\lambda c|)$} &
\egcell{$e_0=\sum\delta w+10^{-13}\sum|w|$; $e_G=e_0+10^{-13}(|S_0|+\sum|\hat\lambda c|)+10^{-300}$} & \egcell{9.1}\\
\egcell{$\tau$, $\delta\tau$} &
\egcell{$\tau_m\ge|t_{mi}|,|t^*_{mi}|$; $\delta\tau_m\ge|t_{mi}-t^*_{mi}|$} &
\egcell{$|t^*|\le|t|+\delta t/4.48\le(|t|+\delta t)(1-\uround)$, as $\delta t\ge9\uround|t|$} & \egcell{none needed} & \egcell{---}\\
\egcell{$\kappa=$ \texttt{ak1}, $k_D$} &
\egcell{$\kappa_i\ge|K^*_{1,i}|$; $|k_{D,i}|(1+10^{-14})\ge|K^*_{D,i}|$} &
\egcell{$k_1=\fl(2\hat\gamma\hat s)$: $3\uround$; $k_D=\fl(2\hat\gamma\fl(\hat s^2))$: $5\uround$; product $\uround$} &
\egcell{factor $1+10^{-14}$} & \egcell{15}\\
\egcell{$S_1$, $e_1$} &
\egcell{$|S_{1,i}-S^*_{1,i}|\le e_1$} &
\egcell{$|wt-w^*t^*|\le|w-w^*||t^*|+|w||t-t^*|$: $\sum(\delta w\,\tau+|w|\delta\tau)$; sum of products (F3): $\gamma_{97}\sum|w|\tau$} &
\egcell{$\sum(\delta w\,\tau+|w|\delta\tau+10^{-13}|w|\tau)$} &
\egcell{9.2 (rounding); data part by the slack of $\delta w$ and $\delta t$}\\
\egcell{$S_2$, $e_2$} &
\egcell{$|S_{2,ij}-S^*_{2,ij}|\le e_2$} &
\egcell{$\sum(\delta w\,\tau^2+2|w|\tau\,\delta\tau)$; $\gamma_{98}\sum|w|\tau^2$; $\tau^2$ by \eqref{eq:eg-P}} &
\egcell{$\sum(\delta w\,\tau^2+2|w|\tau\,\delta\tau+10^{-13}|w|\tau^2)$} & \egcell{9.1}\\
\egcell{$e_3$; $s_3=|T|+e_3$} &
\egcell{$s_{3,ijl}\ge|T^*_{ijl}|$} &
\egcell{$\sum(\delta w\,\tau^3+3|w|\tau^2\delta\tau)$; four-factor products $\gamma_{99}\sum|w|\tau^3$; the sum $|T|+e_3$: $\uround$; $\tau^2,\tau^3$ by \eqref{eq:eg-P}} &
\egcell{$\sum(\delta w\,\tau^3+3|w|\tau^2\delta\tau+10^{-13}|w|\tau^3)$} & \egcell{9.0}\\
\egcell{gradient $\beta$, $e_\beta$ (\texttt{egrad})} &
\egcell{$|\beta_i-\nabla^*_i|\le e_{\beta,i}$} &
\egcell{$|K^*_{1,i}|e_1+4\uround|k_{1,i}S_{1,i}|+\uround|\hat\lambda_i|+\uround|\beta_i|$} &
\egcell{$\kappa_ie_1+10^{-13}(|k_{1,i}S_{1,i}|+|\hat\lambda_i|+|\beta_i|)+10^{-300}$} & \egcell{$>20$}\\
\egcell{Hessian $H$, $e_H$} &
\egcell{$H-e_H\le H^*\le H+e_H$ elementwise} &
\egcell{off-diagonal: $\kappa_i\kappa_je_2+8\uround|H_{ij}|$; diagonal adds $|k_{D,i}|(1+5\uround)e_0+6\uround|k_{D,i}S_0|$ and one sum} &
\egcell{$\kappa_i\kappa_je_2+10^{-13}|H_{ij}|$, diagonal $+|k_{D,i}|(1+10^{-14})e_0+10^{-13}|k_{D,i}S_0|$; then $\times(1+10^{-12})+10^{-300}$} & \egcell{100}\\
\egcell{$Q_L$, $Q_U$} &
\egcell{$Q_L\le\frac12d^\top H^*d\le Q_U$ for $|d|\le r$} &
\egcell{the steps $\fl(H\mp e_H)$: $\uround(|H|+e_H)$, covered by the unused part of $10^{-13}|H|$ and of the factor $1+10^{-12}$; squares, products and sums: $30\uround$ relative} &
\egcell{$\times(1+10^{-12})\mp10^{-300}$} & \egcell{300}\\
\egcell{$\ell$, $q$} &
\egcell{$\ell_m\ge\ell^*_m$; $q\ge q^*$} &
\egcell{$\ell$: $\kappa_i\fl(|t|+\delta t)\ge|K^*_{1,i}||t^*|$, then 10 roundings; $q$: data $\uround$, $\hat s^2\ge s^{*2}$ with $\hat s^2$ by \eqref{eq:eg-P}, 12 roundings} &
\egcell{factor $1+10^{-13}$} & \egcell{90 ($\ell$); 8.8 ($q$)}\\
\egcell{$p=\ell+2q$} &
\egcell{$p_m\ge\ell^*_m+2q^*$} &
\egcell{one rounding} & \egcell{slack of $\ell$ and $q$} & \egcell{---}\\
\egcell{$\eta_m=\fl(\exp(\ell_m)(1+10^{-14}))$ (\texttt{el})} &
\egcell{$\eta_m\ge e^{\ell^*_m}(1-3\uround)$} &
\egcell{\eqref{eq:eg-E} and one product} &
\egcell{factor $1+10^{-14}$; the deficit is absorbed by the factors on $P$} &
\egcell{tolerates $\varepsilon\approx10^{-12}$}\\
\egcell{$R_2$, $R_4$, $P_2$, $P_4$} &
\egcell{$P_2\ge P_2^*$; $P_4\ge\sum_m|w^*_m|R_4(\ell^*_m,q^*)$} &
\egcell{$p^2,p^3,p^4$ by \eqref{eq:eg-P}; at most 10 roundings; $\bar w$: $\uround$; 97-term sum $\gamma_{96}$; together $\le\sci{2.2}{-14}$ relative} &
\egcell{$\times(1+10^{-12})$} & \egcell{45}\\
\egcell{cubic $C_3$ (\texttt{cub}), $U_1$, $U_2$, $P_3$, $P$} &
\egcell{$P\ge\min(P_2^*,P_3^*)$} &
\egcell{$C_3$: accumulation of at most 343 seven-factor products, $\gamma_{348}$, and the division by 6; $U_1\ge\sum_i|K^*_{1,i}S^*_{1,i}|r_i$ since $|S_1|+e_1\ge|S_1^*|$; $U_2\ge q^*(1-\sci{1.4}{-15})$} &
\egcell{$P_3=(C_3+U_1U_2)(1+10^{-12})+P_4$; $P=\min(P_2,P_3)(1+10^{-12})+10^{-300}$; $P=\infty$ if some $\ell_m>600$} & \egcell{25}\\
\egcell{$L_E$} &
\egcell{$L_E\ge\sum_ie_{\beta,i}r_i$} &
\egcell{7 products and the sum: $8\uround$} &
\egcell{$\times(1+10^{-12})$} & \egcell{$>1000$}\\
\egcell{outputs $aL$, $aU$} &
\egcell{$aL\le G-e_G-P+Q_L-L_E$; $aU\ge G+e_G+P+Q_U+L_E$ (exact arithmetic)} &
\egcell{five-term chain and the pad steps: $7\uround\cdot\mathit{tot}$, $\mathit{tot}=|G|+e_G+P+|Q_L|+|Q_U|+L_E$} &
\egcell{$10^{-12}\mathit{tot}+10^{-300}$} & \egcell{$>1000$}\\
\end{longtable}
\endgroup

Some lines need a comment.
In the line for $w$, the exact weight is $w^*=\hat a(1+\delta_a)^{-1}e^{E^*_0}$ and the computed one is $w=\hat a\,e^{E_0}(1+\delta_e)(1+\delta_\times)$ with $|\delta_a|,|\delta_\times|\le\uround$ and $|\delta_e|\le\varepsilon$, so $w^*=w\,e^{E^*_0-E_0}/((1+\delta_a)(1+\delta_e)(1+\delta_\times))$ and $|E^*_0-E_0|\le\delta E$.
This gives the error bound of the table, whose affine form we now derive for $0\le d:=\delta E\le10^{-6}$.
Since $1/k!\le\frac12$ for $k\ge2$, the positive Taylor series gives $(e^d-1)/d=\sum_{k\ge1}d^{k-1}/k!\le1+\frac d2\sum_{j\ge0}d^j=1+\frac d{2(1-d)}\le1+\sci{5.01}{-7}$.
With $D=(1-\uround)^2(1-\varepsilon)\ge1-2\uround-\varepsilon$ we have $1/D-1\le\varepsilon+2.01\uround$ and $(1+\sci{5.01}{-7})/D\le1.000001$, so $e^d/D-1=(e^d-1)/D+(1/D-1)\le1.000001\,d+\varepsilon+2.01\uround$ (the three numerical inequalities were also checked in exact rational arithmetic).
The computed pad is at least $|w|(1.02\,\delta E+1.1\cdot10^{-14})(1-4\uround)$.
It covers the $\delta E$ part with almost 2\% to spare and the rest whenever $\varepsilon\le1.1\cdot10^{-14}(1-4\uround)-2.01\uround$, that is, for $\varepsilon$ up to about $\sci{1.07768}{-14}$; at $\varepsilon=10^{-14}$ the excess of the pad is more than 7\%.
When $E_0<-700$ the program sets $w=0$, and $|w^*|\le|\hat a|(1+\uround)e^{-700+\delta E}<10^{-304}|\hat a|$ is covered by the term $|\hat a|10^{-300}$.
In the lines for $S_1$, $S_2$ and $e_3$, the data part is covered because $\delta w$ exceeds $|w-w^*|$ by almost 2\% and $\delta\tau$ exceeds $|t-t^*|$ by a factor of at least 4.48, while the computed sums fall short of the exact ones by at most $\gamma_{96}$; the same slack absorbs the error $\varepsilon$ of the powers $\tau^2$ and $\tau^3$.
In the cubic, $|L_m(d)^3|$ summed with the weights is bounded through $s_3\ge|T^*|$ and $\kappa\ge|K^*_1|$ over ordered triples, as in \cref{lem:eg-taylor}.
The quantities $\bar w$, $\eta_m$ and $U_2$ may fall short of their exact counterparts by a few units of $\uround$; the factors $1+10^{-12}$ on $P_2$, $P_3$, $P_4$ and $P$ absorb these deficits.

With these lines, for $x=c+d\in X$,
\begin{equation*}
  g^*(x)\ \ge\ G^*+\nabla^*\cdot d+\tfrac12d^\top H^*d-\min(P_2^*,P_3^*)
  \ \ge\ (G-e_G)+\beta\cdot d-L_E+Q_L-P\ \ge\ aL+\beta\cdot d ,
\end{equation*}
because $|(\nabla^*-\beta)\cdot d|\le\sum_ie_{\beta,i}r_i\le L_E$.
The first inequality is \cref{lem:eg-taylor} for the exact data and the second is \cref{lem:eg-affine}.
The upper bound is symmetric.
This proves item~(ii) and \cref{lem:eg-padding}. \qed

The lines that use an exponential or a power tolerate relative errors of $\sci{1.077}{-14}$ (term weights), $\sci{1.954}{-14}$ (natural term ends), about $10^{-13}$ ($\hat s^2$) and $10^{-12}$ ($e^{\ell}$, $p^k$), and at least $\sci{5}{-8}$ ($\tau^2,\tau^3$); all exceed the tolerance $\varepsilon=10^{-14}$ that the audit enforces.
In the auditor's tests NumPy's exponential and powers had relative errors of at most $1.13\,\uround$, about 80 times below $\varepsilon$.

\subsection{The accuracy auditor}\label{app:egrounding-auditor}

The accuracy auditor (the auditor, below) is new code written for this purpose: it checks each library exponential and power value that the enclosure code uses.
It shares no code with the exponential or interval modules of the first code or with the interval certifier of \cref{app:eg-proof}.
It uses only binary64 operations covered by \cref{hyp:fp-ieee} and exact integer and rational arithmetic; no library or mpmath result is trusted.
All its constants are derived when the module is loaded, and every inequality used below is asserted there in exact rational arithmetic; loading fails if one of them does not hold.

\begin{lemma}[exponential enclosure]\label{lem:eg-auditexp}
Assume \cref{hyp:fp-ieee}.
For every float $x\in[-708,709]$ the auditor either reports failure or computes floats $lo\le e^x\le hi$.
\end{lemma}

\begin{proof}
The auditor uses a Cody--Waite argument reduction.
Let $L=\ln2/64$.
The partial sum $s_{400}=\sum_{k\le400}1/(k2^k)$ satisfies $s_{400}\le\ln2\le s_{400}+1/(400\cdot2^{400})$, so $L$ is known to within $10^{-120}$.
Let $L_1$ be $L$ truncated to 36 significant bits and $L_2=\fl(L-L_1)$; then $|L-L_1-L_2|\le D_L=\sci{1.6}{-30}$, an exact bound.
The auditor computes $m=\mathrm{rint}(\fl(x\cdot\tilde L^{-1}))$ with a float approximation $\tilde L^{-1}$ of $1/L$ (any integer $m$ would do), checks $|m|<2^{17}$, so that $p_1=mL_1$ is exact, and computes $s=\fl(x-p_1)$, $p_2=\fl(mL_2)$ and $r=\fl(s-p_2)$.
It checks $|r|\le0.0055$.
Let $r_{\mathrm{true}}=x-mL$.
By (F0), $|r_{\mathrm{true}}-r|\le\uround(|r|+|s|+|p_2|)+|m|D_L+3\eta_0\le D_R=\sci{1.3}{-18}$.

Let $H$ be the value of the degree-7 Taylor polynomial $\sum_{i\le7}c_ir^i$ computed by Horner's rule, with float coefficients $c_i=\fl(1/i!)$.
Unrolling Horner's rule shows that the computed value equals $\sum_ic_ir^i(1+\theta_i)$ with $|\theta_i|\le\gamma_{14}$, since the contribution of each coefficient passes through at most 14 roundings.
Adding the coefficient errors $\sum_i|c_i-1/i!||r|^i$, the Lagrange remainder $|r|^8e^{|r|}/8!$ and $16\eta_0$ for underflow gives $|e^r-H|\le E_H=\sci{1.57}{-15}$; and $H\ge1-0.0055-E_H$, so $|e^r-H|\le\rho_HH$ with $\rho_H=\sci{1.572}{-15}$.

Write $m=64k+j$ with $0\le j<64$.
Then $e^x=2^k\,2^{j/64}e^{r_{\mathrm{true}}}$.
Floats $T^{\mathrm{lo}}_j\le2^{j/64}\le T^{\mathrm{hi}}_j$ are proved by $(T^{\mathrm{lo}}_j)^{64}\le2^j\le(T^{\mathrm{hi}}_j)^{64}$ in integer arithmetic.
The auditor returns
\begin{equation*}
  lo=\fl\bigl(\fl(T^{\mathrm{lo}}_jH)\,C^{\mathrm{lo}}\bigr)\cdot2^k,\qquad hi=\fl\bigl(\fl(T^{\mathrm{hi}}_jH)\,C^{\mathrm{hi}}\bigr)\cdot2^k,
\end{equation*}
where the float constants satisfy $C^{\mathrm{lo}}(1+\uround)^2\le(1-\rho_H)(1-D_R)$ and $C^{\mathrm{hi}}(1-\uround)^2\ge(1+\rho_H)(1+2D_R)$ (checked exactly).
Since $e^{r_{\mathrm{true}}}\ge e^r(1-D_R)\ge H(1-\rho_H)(1-D_R)$ and $e^{r_{\mathrm{true}}}\le H(1+\rho_H)e^{D_R}\le H(1+\rho_H)(1+2D_R)$, the two products bracket $2^{j/64}e^{r_{\mathrm{true}}}$.
The factor $2^k$ is built from its bit pattern.
For $x\ge-708$ every result is normal, because $e^{-708}(1-10^{-14})>2^{-1022}$ (checked exactly), and for $x\le709$ none overflows (checked exactly); so the scaling by $2^k$ is exact and $lo\le e^x\le hi$.
\end{proof}

\paragraph{The exponential check.}
For $-708\le x\le709$ a call $y=\exp(x)$ passes if and only if the reduction checks succeed and
$y\le\fl(lo\cdot F^{\mathrm{up}})$ and $y\ge\fl(hi\cdot F^{\mathrm{dn}})$, with float constants $F^{\mathrm{up}}(1+\uround)\le1+\varepsilon$ and $F^{\mathrm{dn}}(1-\uround)\ge1-\varepsilon$ (checked exactly).
Then $(1-\varepsilon)e^x\le y\le(1+\varepsilon)e^x$ by \cref{lem:eg-auditexp}.
For $x<-708$ the call passes if and only if $0\le y\le2^{-1020}$.
Every other case fails, including NaN, $x>709$ and failed reduction checks.
So a passed call proves \eqref{eq:eg-E} for that call.

\begin{lemma}[power check]\label{lem:eg-auditpow}
Assume \cref{hyp:fp-ieee}.
Let a call $y=x^k$ with $k\in\{2,3,4\}$ pass the auditor's check, and let $x=0$, or $0<x<2^{-250}$, or $2^{-200}\le x\le2^{250}$.
Then \eqref{eq:eg-P} holds for this call.
\end{lemma}

\begin{proof}
For $x\in[2^{-200},2^{250}]$ the auditor computes $z=\fl(x\cdot x)$ for $k=2$, $z=\fl(\fl(x\cdot x)\cdot x)$ for $k=3$ and $z=\fl(\fl(x\cdot x)^2)$ for $k=4$.
No product underflows or overflows, so $|z-x^k|\le G_zx^k$ with $G_z=(1+\uround)^3-1$.
The call passes if and only if $|\fl(y-z)|\le\fl(C_pz)$, where the float $C_p>0$ satisfies $C_p(1+G_z)(1+\uround)+G_z\le\varepsilon$ (checked exactly).
Since $z\ge2^{-801}$ and $C_p>\sci{9}{-15}$, the product $C_pz$ is normal, so $\fl(C_pz)\le C_pz(1+\uround)<z/2$.
If $y<z/2$ or $y>2z$, then by monotone rounding $|\fl(y-z)|\ge z/2>\fl(C_pz)$, and the call fails.
Otherwise $y-z$ is exact by Sterbenz's lemma, and $|y-x^k|\le|y-z|+|z-x^k|\le C_pz(1+\uround)+G_zx^k\le\bigl(C_p(1+G_z)(1+\uround)+G_z\bigr)x^k\le\varepsilon x^k$.
For $x=0$ the call passes only if $y=0$; for $0<x<2^{-250}$ the condition is checked in exact rational arithmetic; negative and NaN arguments and arguments above $2^{250}$ fail.
\end{proof}

The lemma does not cover arguments in $[2^{-250},2^{-200})$, where for $k=4$ the product $C_pz$ can be subnormal and the floating-point comparison needs an extra absolute term.
No power argument of the re-certifier lies there.
The arguments are $\tau$, $p$ and $\hat s$, and on the root boxes $\tau>\sci{2.99}{-16}$, $p$ is $0$ or larger than $10^{-34}$, and $\hat s\in\{1,\fl(0.1)\}$.

\paragraph{Bookkeeping.}
The audit reduces the checks of one call to one flag per box.
The flags are propagated through the recursive splitting of the re-certifier, and a leaf counts as audit-clean only if no call made for it or for any of its pieces failed.
A failure in a piece that was later split also flags the leaf, which is conservative.
The power $\hat s^2$ is computed once per batch of 64 boxes, so a failure there flags every box of the batch.
The certification decisions themselves are not changed by the audit.

\subsection{The audited rerun}\label{app:egrounding-run}

The audited copy of the re-certifier differs from the original only by wrappers that pass every exponential and power result, unchanged, to the auditor; removing the wrappers mechanically recovers the original program line for line, and the copies of the driver and of the margin-recording layer differ from their originals only in marked lines.
The rerun uses the same recorded trees and the same chunking as the original runs, so every 64-box Taylor batch, every LP and every exact test is the same.
The deleted recordings of \inst{eg_int_s} and \inst{eg_disc_s} were regenerated by a deterministic replay of the search, which reproduced the original search logs line for line apart from timings.

A comparison program checked the outputs of the rerun.
For \inst{eg_disc2_s}, the per-leaf results of all 38 chunks (certified or not, the test used and the margin) are identical byte for byte to the saved results of the original run.
For \inst{eg_int_s} and \inst{eg_disc_s}, whose original per-leaf results were not kept, every progress line, the full statistics and the smallest LP margin agree with the logs of the original runs, and the rerun is a complete certificate in its own right.
Both coverage checks of \cref{app:eg-computation} and the domain check passed on the leaf sets of the rerun.
Every call site was exercised; at each of the four per-term sites the number of audited exponentials equals the number of pieces times $28\times97$.

Tests of the auditor are evidence about the implementation; the soundness of the check rests on \cref{lem:eg-auditexp,lem:eg-auditpow}.
Two separately written test programs compared the enclosure of \cref{lem:eg-auditexp} with mpmath at 300 and 400 bits on 173{,}716 arguments in all, including reduction half-points, the ends $-708$ and $709$ and subnormal arguments, and checked that injected errors above $\varepsilon$ and NaN, infinite and negative results are rejected and errors of at most $\frac12\varepsilon$ accepted; end-to-end injections on 512 leaves flagged exactly the leaf that owns the affected piece.

\begin{table}[ht]
\centering
\footnotesize
\setlength{\tabcolsep}{4pt}
\caption{Audited rerun of the leaf re-certification at the targets $\theta^*_I$ (complete run, 2026-10-04).
Pieces: the leaves and the boxes obtained from them by bisection; each piece passed a test of \cref{lem:eg-boxtests} or was bisected (column split).
The margin of a leaf is the smallest certified lower bound on $\Phi$ over its pieces minus $\theta^*_I$ ($+\infty$ for side and Farkas tests); the column shows the smallest margin, rounded down.
No audited value violated its test, and no leaf was flagged.
Times are summed wall times of single-threaded jobs, including the regeneration of the recordings for \inst{eg_int_s} and \inst{eg_disc_s}, on a shared machine.}
\label{tab:eg-audit}
\begin{tabular}{@{}lrrrrrrr@{}}
\toprule
instance & leaves & pieces & split & smallest margin & exponentials & powers & time (s)\\
\midrule
\inst{eg_int_s} & 33{,}385 & 38{,}173 & 2{,}394 & $\sci{7.68}{-11}$ & 414{,}711{,}472 & 829{,}428{,}916 & 465\\
\inst{eg_disc_s} & 86{,}796 & 125{,}382 & 19{,}293 & $\sci{5.40}{-13}$ & 1{,}362{,}150{,}048 & 2{,}724{,}320{,}842 & 1{,}456\\
\inst{eg_disc2_s} & 1{,}114{,}361 & 1{,}769{,}947 & 327{,}793 & $\sci{1.01}{-9}$ & 19{,}228{,}704{,}208 & 38{,}457{,}813{,}736 & 14{,}782\\
\midrule
total & 1{,}234{,}542 & 1{,}933{,}502 & 349{,}480 & & 21{,}005{,}565{,}728 & 42{,}011{,}563{,}494 & 16{,}703\\
\bottomrule
\end{tabular}
\end{table}

\Cref{tab:eg-audit} summarizes the result: 63{,}017{,}129{,}222 exponential and power results were audited, with no violation, and all 1{,}234{,}542 leaves were certified and audit-clean.

\subsection{Trust base of the dual bounds}\label{app:egrounding-trust}

\begin{theorem}[dual bounds of the \inst{eg} instances]\label{thm:eg-trust}
Assume \cref{hyp:fp-ieee}.
Assume that the following programs do what \cref{app:eg} and this section describe, and that the archived outputs of the audited run are their outputs on the stated inputs and targets: the GAMS reader of the re-certifier and the exact comparison of its data with the OSIL data; the enclosure procedures and exact tests of the re-certifier; the auditor; and the coverage, domain and comparison checks.
Then, for each of the three instances $I$, every exactly feasible point of $\model_I$ has objective value at least $\theta^*_I$.
In particular $\vstar(\model_I)$ is at least $6.4531031529331155$ for \inst{eg_int_s}, $5.760539610694993$ for \inst{eg_disc_s} and $5.642100574331458$ for \inst{eg_disc2_s}.
\end{theorem}

\begin{proof}
By the audit (\cref{tab:eg-audit}) and \cref{lem:eg-auditexp,lem:eg-auditpow}, every exponential and power result that the re-certifier used for any piece of any leaf satisfies \eqref{eq:eg-E} or \eqref{eq:eg-P}; all power arguments lie in the range covered by \cref{lem:eg-auditpow}.
Every piece lies in a root box, so by \cref{lem:eg-padding} the floats computed for it satisfy \eqref{eq:eg-rowbounds} for the exact rows.
Each piece passed one of the tests of \cref{lem:eg-boxtests} at $\theta^*_I$, evaluated exactly.
The coverage and domain checks establish the hypothesis of \cref{lem:eg-cover}, which gives $\vstar(\model_I)\ge\theta^*_I$.
The displayed values are at most $\theta^*_I$.
\end{proof}

Every exponential and integer-power result used in the enclosures was checked against a proved enclosure or error test, so no accuracy guarantee for the exponential and power routines of NumPy, the C mathematical library or Intel SVML is assumed; the LP solver only proposes multipliers, whose validity is checked exactly.
Correct implementation includes Python integer and \texttt{Fraction} arithmetic and conversions to binary64, NumPy's basic arithmetic and reductions, comparisons, array indexing, adjacent-float operations and the integer and bit operations of the auditor.
The first code's search and its error analyses, solver optimality claims, sampled accuracy measurements and mpmath are not premises of the theorem.
A separate agent session read the audit code and an earlier written form of the proofs in this section, re-derived the auditor's constants, repeated the comparison and coverage checks, and replayed the audited certification of \inst{eg_int_s} and of part~7 of \inst{eg_disc2_s}, reproducing every saved field except elapsed time; it found four minor issues, which this section corrects.
Nothing is formally verified.

\begin{certbox}[Certificate for \cref{thm:eg-bounds} (dual part: \cref{thm:eg-trust})]
\certitem{Inputs}{OSIL files (SHA-256 prefixes \texttt{a462a01b}, \texttt{5f21b8d7}, \texttt{8849e06b}); GAMS files (\texttt{0fd00330}, \texttt{f38457d6}, \texttt{3ad4d4ea}), equal to the OSIL data as rationals; the recorded leaf boxes of the first code's search (1{,}234{,}542 leaves); the points of \cref{tab:eg-points}.}
\certitem{Computation}{dual: per leaf, natural and Taylor-model enclosures in binary64 with the paddings of \cref{lem:eg-padding} \arith{F}; box tests at $\theta^*_I$ in exact rational arithmetic \arith{E}; bisection of failing boxes; every exponential and power result checked by the auditor (\cref{lem:eg-auditexp,lem:eg-auditpow}); coverage proved exactly from the leaf set alone. Primal: all 28 rows in library-free dyadic interval arithmetic \arith{E}; bounds and integrality exact.}
\certitem{Data reading}{dual: term data as correctly rounded doubles, covered by (F1) in \cref{lem:eg-padding}; $c_k$, side-row bounds, variable bounds and $\theta^*_I$ exact. Primal: decimals enclosed outward in dyadic intervals by the check's own OSIL reader.}
\certitem{Implementations}{dual: the second code, the leaf re-certifier, on all leaves carries the displayed bounds; the first code's search in interval mode certifies \inst{eg_int_s} and \inst{eg_disc_s} completely and part~1 of \inst{eg_disc2_s} under a side condition (\cref{tab:eg-matrix}). Primal: the dyadic check, rerun and read in full by a separate agent session and confirmed by decimal and mpmath interval evaluations with two further readers (\cref{app:eg-primal}).}
\certitem{Evidence level}{dual: \evid{hand} for the lemmas and \evid{rerun} for the leaf certificates (the leaf partition is stored and its coverage is checked exactly; only the per-leaf certificates are recomputed); primal: \evid{stored}.}
\certitem{Replay}{dual, audited re-certification including the regeneration of two recordings (summed single-threaded wall times, \cref{tab:eg-audit}): Tier~1 for \inst{eg_int_s} (465~s), Tier~2 for \inst{eg_disc_s} (1{,}456~s), Tier~3 for \inst{eg_disc2_s} (14{,}782~s); 37 minutes for all three with eight parallel jobs. Primal: Tier~1; about 17~s for \inst{eg_int_s}.}
\end{certbox}
